\documentclass{article}

% NeurIPS 2025 style — download neurips_2025.sty from the official submission kit
% and place it alongside this file.
\usepackage{neurips_2025}

% ---------------------------------------------------------------
% Standard packages
% ---------------------------------------------------------------
\usepackage[utf8]{inputenc}
\usepackage[T1]{fontenc}
\usepackage{hyperref}
\usepackage{url}
\usepackage{booktabs}
\usepackage{amsfonts}
\usepackage{amsmath}
\usepackage{nicefrac}
\usepackage{microtype}
\usepackage{xcolor}
\usepackage{graphicx}
\usepackage{multirow}
\usepackage{array}
\usepackage{enumitem}

% ---------------------------------------------------------------
% Outline helpers — remove before submission
% ---------------------------------------------------------------
\newcommand{\outline}[1]{\textcolor{gray}{\textit{[#1]}}}
\newcommand{\todo}[1]{\textcolor{red}{\textbf{[TODO: #1]}}}
\newcommand{\N}{$N$}

% ---------------------------------------------------------------
% Paper metadata
% ---------------------------------------------------------------
\title{ReproScore: An Agentic Framework for Automated\\
       Reproducibility Assessment of Scientific Publications}

\author{%
  \textbf{Author 1} \\
  Affiliation \\
  \texttt{email@institution.edu}
  \And
  \textbf{Author 2} \\
  Affiliation \\
  \texttt{email@institution.edu}
}

% ---------------------------------------------------------------
\begin{document}

\maketitle

% ==============================================================
\begin{abstract}
% Target: ~250 words
\outline{%
  Open with the reproducibility crisis (quantify the scope: 1.82\% of
  transportation simulation studies provide reproducible repositories;
  broader estimates suggest fewer than one-third of computational studies
  can be independently reproduced).
  %
  State the core problem: existing reproducibility audits rely on manual
  human effort and binary pass/fail verdicts that do not scale and provide
  no actionable guidance to authors.
  %
  Introduce \textbf{ReproScore}: an agentic pipeline that automatically
  assigns a continuous, multi-dimensional Reproducibility Index (0--100)
  to any computational paper given its PDF or DOI.
  %
  Briefly describe the five agents (Analyst, Retriever, Builder, Runner,
  Auditor) and the five scoring dimensions (Artefact Availability,
  Environment Reproducibility, Execution Success, Output Fidelity,
  Documentation Quality).
  %
  Validation sentence: we validate ReproScore against \N{} papers from
  ReScience~C, where independent human reviewers rendered reproduction
  verdicts; report top-line agreement metric (e.g., band accuracy,
  Spearman $\rho$).
  %
  Close with broader impact: journals, preprint servers, and funding
  agencies can integrate ReproScore into submission and review workflows
  without additional human overhead.
}
\end{abstract}

% ==============================================================
\section{Introduction}
\label{sec:intro}

\outline{%
  \textbf{Hook (1--2 paragraphs).}
  The reproducibility crisis in computational science.
  Cite: meta-analysis finding that only 1.82\% of transportation
  simulation papers (2000--2024, $n = 46{,}015$) provide working
  repositories~\cite{riehl2025}; 64\% of researchers cite time
  constraints as the primary barrier to sharing materials.
  Broader estimates: fewer than 40\% of ML results reproduce under
  independent re-implementation.
  Cost: wasted research effort, erosion of public trust.
}

\outline{%
  \textbf{Limitations of current solutions (1 paragraph).}
  Manual reproduction efforts (ReScience~C~\cite{rescience}, ML
  Reproducibility Challenge) deliver high-quality verdicts but cannot
  scale to the thousands of papers published weekly.
  Static checklists and author-reported artifact badges (ACM, Nature)
  are unverified and binary.
  No existing system automatically \emph{executes} a paper's code and
  compares outputs to published results.
}

\outline{%
  \textbf{Our approach (1 paragraph).}
  ReproScore addresses the gap: given any computational paper, it
  automatically (i) extracts methodology, datasets, and artifacts,
  (ii) retrieves and builds the computational environment,
  (iii) executes the pipeline in a sandbox, and (iv) compares outputs
  to published figures and tables using quantitative thresholds,
  yielding a reproducibility score that correlates with human verdicts.
}

\outline{%
  \textbf{Contributions (bullet list).}
  \begin{enumerate}[label=\arabic*.]
    \item \textbf{Reproducibility Index.} A 5-dimensional, continuous
          scoring scheme (0--100) with validated per-dimension thresholds
          (SSIM $\geq$ 0.85 for figures, $\leq$1\% per-cell error for
          tables, etc.).
    \item \textbf{ReproScore pipeline.} An open-source, domain-agnostic
          agentic framework that assigns the index automatically from a
          paper PDF or DOI, with human-in-the-loop gates for ambiguous cases.
    \item \textbf{Empirical validation.} Agreement analysis between
          ReproScore and human reviewer verdicts on \N{} ReScience~C papers,
          including ablation over scoring dimensions.
    \item \textbf{Case studies.} Three end-to-end traces in the
          transportation domain illustrating graceful degradation when
          artifacts are unavailable.
  \end{enumerate}
}

% ==============================================================
\section{Related Work}
\label{sec:related}

\outline{%
  Keep to 1--1.5 pages. Organise into four paragraphs below.
  End each paragraph with a sentence positioning ReproScore against that
  line of work.
}

\paragraph{Manual reproducibility initiatives.}
\outline{%
  ReScience~C~\cite{rescience}: peer-reviewed journal for explicit
  replications.
  ML Reproducibility Challenge: community effort; covers only ML papers;
  not scalable; no continuous score.
  Papers With Code leaderboards: track results but do not verify them.
  \textbf{Gap:} all are manual; none produce a continuous score.
}

\paragraph{Artifact badging and reporting checklists.}
\outline{%
  ACM artifact review (Functional / Reusable / Available badges).
  Nature and Science structured reporting checklists.
  NeurIPS reproducibility checklist.
  Limitations: author-reported, unverified, binary.
  \textbf{Gap:} badges require a human reviewer; checklists do not test
  execution.
}

\paragraph{Automated code and metadata analysis.}
\outline{%
  Static analysis tools (CodeCarbon, dependency scanners).
  Automated metadata extraction from papers (CERMINE, GROBID, SciSpacy).
  Recent LLM-based extraction~\cite{todo_llm_extraction}: high recall but
  hallucination risk for version numbers and dataset URLs.
  \textbf{Gap:} no existing tool executes code and validates outputs.
}

\paragraph{Source-grounded information extraction.}
\outline{%
  Retrieval-augmented generation (RAG) reduces hallucination but does not
  guarantee answers are grounded in the \emph{specific} source document.
  NotebookLM-style closed-corpus grounding: answers are restricted to
  uploaded documents, eliminating cross-document contamination.
  We use this property in Agent~1 (Section~\ref{sec:analyst}).
}

% ==============================================================
\section{The Reproducibility Index}
\label{sec:index}

\outline{%
  Self-contained section (~1 page). Readers should be able to apply the
  index manually after reading this section.
}

\subsection{Dimensions}
\label{sec:dimensions}

\outline{%
  Introduce the five dimensions as orthogonal axes of reproducibility.
  Justify orthogonality: a paper can have perfect documentation but
  broken code (low Execution Success, high Documentation Quality).
  Present as a table:
}

\begin{table}[h]
  \centering
  \caption{The five dimensions of the Reproducibility Index (20 points each).}
  \label{tab:dimensions}
  \begin{tabular}{clp{7cm}}
    \toprule
    \textbf{Pts} & \textbf{Dimension} & \textbf{What it measures} \\
    \midrule
    20 & Artefact Availability       & Code, data, and models are publicly retrievable at the claimed locations. \\
    20 & Environment Reproducibility & All dependencies resolve without manual intervention; environment builds cleanly. \\
    20 & Execution Success           & The computational pipeline runs to completion without errors. \\
    20 & Output Fidelity             & Reproduced outputs match published figures and tables within defined tolerances. \\
    20 & Documentation Quality       & README and inline comments are sufficient for independent reproduction. \\
    \bottomrule
  \end{tabular}
\end{table}

\subsection{Quantitative Thresholds}
\label{sec:thresholds}

\outline{%
  For each output modality, state the threshold and justify it.
  \begin{itemize}
    \item \textbf{Figures:} Structural Similarity Index (SSIM)~\cite{wang2004ssim}
          $\geq 0.85$; perceptual hash distance $\leq 10$.
          Justify: SSIM 0.85 corresponds to visually near-identical images
          in standard benchmarks.
    \item \textbf{Tables:} Per-cell relative error $\leq 1\%$;
          report percentage of cells within tolerance.
    \item \textbf{Statistical results:} $p$-values within one order of
          magnitude; regression coefficients within 5\% relative error;
          $R^2$ within 0.02 absolute.
    \item \textbf{Stochastic outputs:} Pipeline run 5$\times$ with different
          seeds; distributions compared (KS test, $p > 0.05$).
  \end{itemize}
  Note that thresholds are calibrated on the ReScience~C validation set
  (Section~\ref{sec:validation}) and may be tightened in domain-specific plugins.
}

\subsection{Score Interpretation and Mapping to Existing Schemes}
\label{sec:mapping}

\outline{%
  \begin{itemize}
    \item Present the five bands: 90--100 Fully Reproducible,
          70--89 Largely Reproducible, 40--69 Partially Reproducible,
          10--39 Minimally Reproducible, 0--9 Not Reproducible.
    \item Map bands to ACM artifact badges and FAIR data principles
          (Findable, Accessible, Interoperable, Reusable)~\cite{wilkinson2016}.
    \item Note that the mapping allows ReproScore to serve as a
          \emph{verified} replacement for self-reported badges.
  \end{itemize}
}

% ==============================================================
\section{The ReproScore Pipeline}
\label{sec:pipeline}

\outline{%
  ~2 pages. Lead with the system overview figure (Mermaid/TikZ diagram
  of the state machine and five agents). Then one subsection per agent.
  Emphasise design decisions: source-grounding, containerisation,
  append-only ledger.
}

\subsection{System Overview}
\label{sec:overview}

\outline{%
  State machine: INTAKE $\to$ ANALYSE $\to$ RETRIEVE $\to$ BUILD
  $\to$ EXECUTE $\to$ VALIDATE $\to$ REPORT.
  Each transition logged to the append-only JSON-lines audit ledger.
  Graceful degradation: any blocking failure transitions to
  PARTIAL\_REPORT rather than silent failure.
  \textbf{[Figure: architecture diagram]}
}

\subsection{Agent 1 — Paper Analyst}
\label{sec:analyst}

\outline{%
  \textbf{Input:} PDF or DOI.
  \textbf{Mechanism:} DOI resolved via CrossRef API for metadata.
  Paper uploaded to a closed-corpus LLM (Google NotebookLM via MCP);
  a 5-query sequence targets each \texttt{PaperProfile} field:
  (1) methodology steps and tools, (2) datasets and licenses,
  (3) artifact locations, (4) figures/tables to reproduce,
  (5) hyperparameters and hardware requirements.
  Source-grounding eliminates hallucinated tool versions or dataset URLs;
  confidence scored from citation density in responses.
  \textbf{Fallback:} Gemini API with PDF upload (lower source-grounding guarantee).
  \textbf{Output:} \texttt{PaperProfile} JSON (Pydantic schema).
}

\subsection{Agent 2 — Artefact Retriever}
\label{sec:retriever}

\outline{%
  Resolves repository and dataset links from \texttt{PaperProfile}.
  Fallback chain for dead links: Wayback Machine $\to$ GitHub search
  $\to$ author profile $\to$ Papers With Code $\to$ Zenodo/Figshare.
  Handles ``available upon request'' data: logs as UNAVAILABLE\_GATED;
  optionally generates synthetic data from published summary statistics.
  \textbf{Output:} \texttt{ArtefactInventory} JSON.
}

\subsection{Agent 3 — Environment Builder}
\label{sec:builder}

\outline{%
  Reads dependency manifests (\texttt{requirements.txt},
  \texttt{environment.yml}, \texttt{renv.lock}, \texttt{Project.toml}).
  Version resolution strategy: pin $\to$ range $\to$ latest (progressive
  relaxation on conflict).
  Builds a Docker image for isolation; image hash logged for provenance.
  Flags proprietary software dependencies (e.g., MATLAB, VISSIM) and
  suggests open-source alternatives.
  \textbf{Output:} Docker image tag + build log.
}

\subsection{Agent 4 — Execution Runner}
\label{sec:runner}

\outline{%
  Executes the pipeline inside the Docker container with resource limits
  (CPU, memory, wall-clock time).
  Human-in-the-loop gate triggered if runtime exceeds 1 hour or GPU is
  required (COMPUTE\_APPROVAL).
  Captures stdout/stderr, exit codes, and all generated artifacts.
  Stochastic components run 5$\times$ with different random seeds.
  \textbf{Output:} Execution logs + reproduced artifact files.
}

\subsection{Agent 5 — Validation Auditor}
\label{sec:auditor}

\outline{%
  Compares reproduced artifacts to targets extracted by Agent~1.
  Per-modality comparison (thresholds from Section~\ref{sec:thresholds}).
  Computes per-dimension scores (Table~\ref{tab:dimensions}) and
  aggregates into the Reproducibility Index.
  Generates human-readable recommendations for each failed check.
  \textbf{Output:} \texttt{ReproducibilityReport} JSON + PDF report.
}

\subsection{Human-in-the-Loop Gates}
\label{sec:hitl}

\outline{%
  Four defined gates prevent fully autonomous action on ambiguous inputs:
  \begin{enumerate}[label=\arabic*.]
    \item \textbf{LICENCE\_REVIEW} — non-OSI license or missing license.
    \item \textbf{CODE\_SAFETY} — untrusted code with network calls or
          file writes outside sandbox.
    \item \textbf{COMPUTE\_APPROVAL} — runtime $>$ 1 hour or GPU required.
    \item \textbf{AMBIGUITY\_RESOLUTION} — multiple candidate datasets
          or repos; paper description insufficient to disambiguate.
  \end{enumerate}
  Justify: fully autonomous execution of arbitrary code introduces legal
  and security risks; gates bound the blast radius.
}

\subsection{Audit Trail}
\label{sec:ledger}

\outline{%
  Every agent action appended to an immutable JSON-lines ledger:
  \texttt{\{timestamp, phase, action, data\}}.
  Enables full provenance reconstruction; supports re-running from any
  checkpoint after a gate resolution.
  Ledger is a first-class output of ReproScore and can itself be archived
  for meta-reproducibility (reproducing the reproduction).
}

% ==============================================================
\section{Validation Against ReScience~C}
\label{sec:validation}

\outline{%
  Core empirical section. Target: ~1.5 pages. This is what makes the
  paper publishable: we have ground truth from human experts.
}

\subsection{ReScience~C as Ground Truth}
\label{sec:rescience}

\outline{%
  Describe ReScience~C: peer-reviewed journal for explicit replications of
  computational studies; each submission attempts to reproduce a prior
  paper and is reviewed by independent experts.
  Verdicts map to three bands: Fully Reproduced, Partially Reproduced,
  Not Reproduced.
  Justify as ideal ground truth: independent of authors, transparent
  review, documented failure modes, multi-domain.
  Acknowledge limitation: reviewer expertise and effort vary; we treat
  verdicts as noisy but reliable labels.
}

\subsection{Dataset Construction}
\label{sec:dataset}

\outline{%
  Inclusion criteria: (i) computational paper with claimed code/data,
  (ii) original paper PDF accessible, (iii) ReScience~C article provides
  a clear verdict.
  Exclusion: hardware-specific papers (FPGAs, custom chips) where
  execution is infeasible without specialized hardware.
  Annotation mapping: Fully Reproduced $\to$ expected score $\in [70, 100]$;
  Partially Reproduced $\to [40, 69]$; Not Reproduced $\to [0, 39]$.
  \textbf{[Table: dataset statistics — N papers, domain breakdown, year range,
  verdict distribution]}
  \todo{Fill in N and statistics after dataset construction.}
}

\subsection{Experimental Protocol}
\label{sec:protocol}

\outline{%
  For each paper in the validation set:
  \begin{enumerate}[label=\arabic*.]
    \item Run the full ReproScore pipeline (or log the gate that halted it).
    \item Record per-dimension scores and aggregate Reproducibility Index.
    \item Classify into the three verdict bands using the mapping in
          Section~\ref{sec:mapping}.
    \item Compare to the ReScience~C human verdict.
  \end{enumerate}
  Metrics: band-level accuracy (exact match), Spearman rank correlation
  between aggregate score and ordinal human verdict, confusion matrix,
  per-dimension Pearson correlation with verdict.
  Hardware and runtime budget per paper: \todo{specify}.
}

\subsection{Results}
\label{sec:results}

\outline{%
  \textbf{[Table: confusion matrix — ReproScore band vs.\ human verdict]}\\
  \textbf{[Table: per-metric results — band accuracy, Spearman $\rho$,
  per-dimension correlation]}\\
  Key findings to discuss:
  \begin{itemize}
    \item Overall band accuracy and $\rho$; compare to a majority-class baseline.
    \item Which dimension is the strongest predictor of human verdict?
          (Hypothesis: Output Fidelity, but Artefact Availability may
          dominate due to many ``not available'' cases.)
    \item Error analysis: where does ReproScore disagree with human reviewers
          and why? (e.g., human reviewer had domain expertise to fix an
          undocumented step; pipeline could not).
  \end{itemize}
}

\subsection{Ablation Studies}
\label{sec:ablation}

\outline{%
  \begin{enumerate}[label=\arabic*.]
    \item \textbf{Dimension ablation.} Remove each dimension in turn;
          measure drop in band accuracy. Quantifies each dimension's
          marginal contribution.
    \item \textbf{Source-grounding ablation.} Replace Agent~1 (NotebookLM)
          with a naive prompted LLM (GPT-4o / Gemini); measure downstream
          accuracy loss due to hallucinated metadata.
    \item \textbf{Threshold sensitivity.} Vary SSIM threshold from 0.70
          to 0.95; report accuracy vs.\ threshold curve to justify the
          chosen 0.85 cutoff.
  \end{enumerate}
}

% ==============================================================
\section{Case Studies}
\label{sec:cases}

\outline{%
  Three end-to-end traces from the transportation domain.
  Goal: qualitative illustration of how ReproScore handles varying levels
  of artifact availability; not the primary validation (that is
  Section~\ref{sec:validation}).
  Present as a summary table + one paragraph per paper.
  Move detailed traces to Appendix~\ref{app:traces} if space is tight.
}

\begin{table}[h]
  \centering
  \caption{ReproScore results for three transportation-domain case studies.}
  \label{tab:cases}
  \begin{tabular}{lcccccc}
    \toprule
    \textbf{Paper} & \textbf{AA} & \textbf{ER} & \textbf{ES} & \textbf{OF} & \textbf{DQ} & \textbf{Total} \\
    \midrule
    Riehl et al.\ (2025)~\cite{riehl2025}  & \todo{} & \todo{} & \todo{} & \todo{} & \todo{} & \todo{} \\
    JRC OBFCM emissions study               & \todo{} & \todo{} & \todo{} & \todo{} & \todo{} & \todo{} \\
    DRL ramp metering study                 & \todo{} & \todo{} & \todo{} & \todo{} & \todo{} & \todo{} \\
    \bottomrule
  \end{tabular}
  \smallskip

  \footnotesize AA = Artefact Availability, ER = Environment Reproducibility,
  ES = Execution Success, OF = Output Fidelity, DQ = Documentation Quality.
  All scores out of 20.
\end{table}

\paragraph{Riehl et al.\ (2025) — Meta-analysis of 46,015 transportation papers.}
\outline{%
  Open repository; Python 3.9; 7 pipeline steps; all figures and tables
  reproducible from raw data.
  Expected high score ($\geq 70$).
  Key insight: documentation quality penalised by absent data dictionary
  for one intermediate CSV.
}

\paragraph{JRC OBFCM EU Emissions Study.}
\outline{%
  7.7 million vehicle records; multi-step regression.
  Data publicly available from EU open data portal.
  Tests retrieval and environment-building agents at scale.
  Expected moderate-to-high score; flag any version conflict in
  statistical packages.
}

\paragraph{DRL Ramp Metering Study.}
\outline{%
  ``Data available upon request''; SUMO simulation required.
  Artefact Availability score: near zero (no public repo, no public data).
  Illustrates graceful degradation: pipeline produces a PARTIAL\_REPORT
  with score = \todo{} and recommendation to authors.
  Execution Success and Output Fidelity cannot be evaluated; documented
  as UNAVAILABLE\_GATED in the audit ledger.
}

% ==============================================================
\section{Discussion}
\label{sec:discussion}

\subsection{Limitations}
\label{sec:limitations}

\outline{%
  \begin{itemize}
    \item \textbf{Execution cost.} Running each paper's code at scale
          requires non-trivial compute; mitigated by the COMPUTE\_APPROVAL
          gate and partial scoring when execution is infeasible.
    \item \textbf{Proprietary dependencies.} Papers relying on MATLAB,
          VISSIM, or licensed datasets cannot be fully scored;
          ReproScore flags these and scores available dimensions only.
    \item \textbf{Extraction noise.} Agent~1 accuracy bounds all
          downstream agents; ablation (Section~\ref{sec:ablation})
          quantifies this.
    \item \textbf{ReScience~C coverage.} Validation set is skewed toward
          ML and computational biology; transportation and physical
          simulation papers are underrepresented.
    \item \textbf{Threshold calibration.} Current thresholds are
          calibrated on the validation set; they may not transfer to
          domains with inherently noisier outputs (e.g., climate simulation).
  \end{itemize}
}

\subsection{Ethical Considerations}
\label{sec:ethics}

\outline{%
  \begin{itemize}
    \item \textbf{Code safety.} Execution sandboxed in Docker with no
          network access and restricted filesystem writes; CODE\_SAFETY
          gate provides a human review layer for flagged patterns.
    \item \textbf{Data privacy.} ReproScore never redistributes gated or
          proprietary data; UNAVAILABLE\_GATED papers are scored on
          available dimensions only.
    \item \textbf{Framing.} Scores are descriptive and author-facing, not
          punitive; the report includes actionable recommendations rather
          than a pass/fail verdict.
    \item \textbf{Potential for misuse.} A low score could be weaponised
          in peer review; we recommend scores be shared with authors before
          reviewers and that journals adopt a minimum improvement period.
  \end{itemize}
}

\subsection{Broader Impact and Future Work}
\label{sec:impact}

\outline{%
  \begin{itemize}
    \item \textbf{Journal integration.} Journals could require
          ReproScore $\geq 70$ at submission, shifting the burden of
          proof from reviewers to authors at submission time.
    \item \textbf{Preprint servers.} arXiv/bioRxiv could display
          ReproScore badges alongside papers, providing readers with
          verified reproducibility signals.
    \item \textbf{Funding audits.} Agencies could run ReproScore as a
          post-award check on funded computational outputs.
    \item \textbf{Future work:} domain-specific plugins (climate science,
          genomics); MCP migration for distributed, language-agnostic
          execution; community-contributed threshold calibration datasets;
          longitudinal tracking of a paper's score as authors improve
          their repositories.
  \end{itemize}
}

% ==============================================================
\section{Conclusion}
\label{sec:conclusion}

\outline{%
  Three paragraphs.
  \textbf{(1)} Restate the problem: reproducibility audits at scale
  require automation; no prior work executes papers and validates outputs.
  \textbf{(2)} Summarise ReproScore: five agents, five scoring dimensions,
  human-in-the-loop gates, validated against \N{} ReScience~C papers
  (top-line metric).
  \textbf{(3)} Call to action: open-source release; invite the NeurIPS
  community to contribute papers to the validation benchmark and to
  integrate ReproScore into their own submission workflows.
}

% ==============================================================
\section*{Acknowledgements}
\outline{%
  Funding sources, compute credits, and any human annotators who verified
  ReScience~C verdicts.
}

% ==============================================================
\bibliographystyle{plainnat}
\bibliography{references}

% ==============================================================
\appendix

\section{PaperProfile and ReproducibilityReport Schemas}
\label{app:schemas}

\outline{%
  Full Pydantic JSON schemas for \texttt{PaperProfile} and
  \texttt{ReproducibilityReport}.
  Include: field names, types, allowed values for enums
  (\texttt{DataAvailability}, \texttt{FigureType}).
}

\section{Analyst Query Prompts}
\label{app:prompts}

\outline{%
  Verbatim text of all five versioned queries sent to the closed-corpus
  LLM (from \texttt{prompts/analyst\_queries\_v1.md}).
  Include the fallback prompt (\texttt{analyst\_fallback\_v1.md}).
}

\section{ReScience~C Dataset Construction}
\label{app:dataset}

\outline{%
  Detailed inclusion/exclusion criteria.
  Inter-rater reliability: if multiple ReScience~C reviewers rendered
  different verdicts on the same paper, describe the resolution protocol.
  Full paper list with DOIs, verdicts, and domain labels.
}

\section{Full Pipeline Traces for Case Studies}
\label{app:traces}

\outline{%
  For each of the three case studies (Section~\ref{sec:cases}):
  per-agent timing, full audit ledger excerpt, per-figure SSIM scores,
  and complete \texttt{ReproducibilityReport} JSON.
}

% ==============================================================
\end{document}
